% Part: computability
% Chapter: computability-theory
% Section: prop-reduce

\documentclass[../../../include/open-logic-section]{subfiles}

\begin{document}

\olfileid[fr]{cmp}{thy}{ppr}
\olsection{Propriétés de la réductibilité}

Nous écrivons $A\leq_m B$ lorsque $A$ se réduit
à~$B$. Cette notation suggère que $A$ n'est
« pas plus difficile » que~$B$, que $B$ est au
moins aussi difficile que~$A$, et que $\le_m$
classe les ensembles ou problèmes de décision
selon leur difficulté, à la manière de l'ordre
$\le$ sur les nombres. Les deux propositions
suivantes étayent cette intuition. La première
établit la transitivité de $\le_m$.\footnote{La relation est un préordre : elle est réflexive et transitive, mais deux ensembles distincts peuvent se réduire l'un à l'autre. L'analogie avec l'ordre numérique exprime ici une comparaison de difficulté.}

\begin{prop}
\ollabel{prop:trans-red}
Si $A\leq_m B$ et $B\leq_m C$, alors $A\leq_m C$.
\end{prop}

\begin{proof}
En composant une réduction de $A$ à~$B$ et une
réduction de $B$ à~$C$, on obtient une réduction
de $A$ à~$C$.
\end{proof}

\begin{prob}
Démontrer la \olref{prop:trans-red} en montrant
que si $f$ réduit $A$ à~$B$ et $g$ réduit $B$
à~$C$ au sens many-one, alors $\comp{f}{g}$
réduit $A$ à~$C$ au même sens.
\end{prob}

\begin{prop}
\ollabel{prop:reduce}
Soient $A$ et~$B$ deux ensembles tels que
$A\leq_m B$.
\begin{enumerate}
\item Si $B$ est effectivement énumérable,
  $A$ l'est aussi.
\item Si $B$ est calculable, $A$ l'est aussi.
\end{enumerate}
\end{prop}

\begin{proof}
Soit $f$ une réduction many-one de $A$ à~$B$.
Pour la première assertion, si $B$ est le
domaine d'une fonction calculable partielle~$g$,
alors $A$ est le domaine de $\comp{f}{g}$ :\footnote{Correction éditoriale : la source mentionne seulement une fonction partielle $g$. Pour obtenir un domaine effectivement énumérable, nous précisons que $g$ est calculable partielle.}
\begin{align*}
x \in A & \text{ ssi } f(x) \in B \\
& \text{ ssi }  g(f(x)) \fdefined.
\end{align*}

Pour la seconde assertion, rappelons que si $B$
est calculable, $B$ et~$\Complement{B}$ sont
effectivement énumérables
(\olref[cmp]{thm:ce-comp}). La fonction~$f$
réduit également $\Complement{A}$ à
$\Complement{B}$ au sens many-one. Par la
première assertion, $A$ et~$\Complement{A}$
sont donc effectivement énumérables. Ainsi,
$A$ est calculable par le \olref[cmp]{thm:ce-comp}.
(On peut aussi vérifier que
$\Char{A}=\comp{f}{\Char{B}}$; si
$\Char{B}$ est calculable, $\Char{A}$ l'est
également.)
\end{proof}

\begin{prob}
Supposons que $f$ soit une réduction many-one
de $A$ à~$B$. Montrer que $f$ réduit aussi
$\Complement{A}$ à $\Complement{B}$.
\end{prob}

\begin{prob}
Montrer que si $f\colon\Nat\to\Nat$ est une
réduction many-one de $A$ à~$B$, alors
$\Char{A}=\comp{f}{\Char{B}}$.\footnote{Correction éditoriale : la source donne $f\colon A\to B$. Une réduction many-one est définie sur tous les entiers, y compris hors de $A$ ; son type est $f\colon\Nat\to\Nat$.}
\end{prob}

\begin{digress}
Une notion plus générale, la \emph{réductibilité
de Turing}, est utile dans d'autres contextes,
notamment pour prouver l'indécidabilité. D'après le
\olref[cmp]{cor:comp-k}, le complémentaire
de~$K_0$ ne se réduit pas à~$K_0$ au sens
many-one, puisqu'il n'est pas effectivement
énumérable. Pourtant, si l'on savait répondre
aux questions concernant~$K_0$, on saurait
aussi répondre à celles concernant son
complémentaire. On dit que $A$ est \emph{Turing
réductible} à~$B$ si une procédure calculable
peut décider l'appartenance à~$A$ en posant des
questions sur l'appartenance à~$B$. Cette notion
est plus souple que la réduction many-one :
cette dernière n'autorise qu'une seule question
sur~$B$, et la réponse « oui » ou « non » doit
être transmise telle quelle pour~$A$.
Si $A$ est Turing réductible à~$B$ et si $B$
est calculable, alors $A$ est calculable.
L'énoncé analogue échoue pour l'énumérabilité
effective : le complémentaire de~$K_0$ est
Turing réductible à~$K_0$ sans être
effectivement énumérable.

Il faut bien distinguer ces diverses notions
de réduction. Dans ce chapitre, nous n'utiliserons
que la réduction many-one. Les deux types évoqués
ci-dessus ont aussi des analogues en complexité,
où les machines de Turing doivent s'exécuter en
temps polynomial : la version de la réduction
many-one s'appelle \emph{réduction de Karp},
et celle de la réduction de Turing,
\emph{réduction de Cook}.
\end{digress}

\end{document}
